\documentclass[a4paper,12pt]{article}

\usepackage[utf8]{inputenc}
\usepackage[T1]{fontenc}
\usepackage[hungarian]{babel}



\usepackage[a4paper,margin=2.5cm]{geometry}



\usepackage{array}
\usepackage{tabularx}
\usepackage{longtable}
\usepackage{multirow}
\usepackage{makecell}


\usepackage{enumitem}

\usepackage{booktabs}

\setlength{\parindent}{0pt}
\setlength{\parskip}{4pt}

\newcommand{\EvaluationTable}[7]{

\renewcommand{\arraystretch}{1.35}

{\fontsize{9pt}{11pt}\selectfont

\begin{tabularx}{\textwidth}{|
>{\raggedright\arraybackslash}p{4.2cm}|
>{\centering\arraybackslash}X|
>{\centering\arraybackslash}X|
>{\centering\arraybackslash}X|
>{\centering\arraybackslash}X|
>{\centering\arraybackslash}p{1.8cm}|}

\hline

\multirow{2}{=}{\centering\textbf{Értékelési\\szempont}}
&
\multicolumn{4}{c|}{\textbf{Adható pontszám}}
&
\multirow{2}{=}{\centering\textbf{Elért\\pontszám}}
\\

\cline{2-5}

&
\textbf{0}
&
\textbf{1--2}
&
\textbf{3--4}
&
\textbf{5}
&
\\

\hline

#1
&
#2
&
#3
&
#4
&
#5
&
\textbf{#6}

\\
\hline

\end{tabularx}

\vspace{3mm}

\textbf{Rövid indoklás (opcionális):}

\vspace{2mm}

#7

\vspace{8mm}

}
}

\begin{document}


\begin{center}

{\Large\textbf{ESZTERHÁZY KÁROLY KATOLIKUS EGYETEM}}

\vspace{0.8cm}

{\LARGE\textbf{SZAKDOLGOZAT BÍRÁLATI LAP}}

\end{center}

\vspace{1cm}

\begin{tabularx}{\textwidth}{>{\bfseries}l X}

A hallgató neve: &
sadads \\[0.4cm]

A hallgató Neptun kódja: &
sadasd \\[0.4cm]

A hallgató szakja: &
Programtervező\_informatikus \\[0.4cm]

A szakdolgozat címe: &
asdasd

\end{tabularx}

\vspace{1cm}

\EvaluationTable
{ 1. A szakdolgozat szerkezeti felépítése, tartalmi tagolása, alaki megjelenése. }
{ A dolgozat nem felel meg az EKKE szakdolgozati követelményeinek. }
{ A követelmények-nek lényegében megfelel, de nagyobb hiányosságokkal. }
{ A követelmé-nyeknek lényegében megfelel, kisebb hiányosságokkal. }
{ Kifogástalan szerkezet, tartalmi tagolás, szép kivitel. }
{ 1 }
{  }

\EvaluationTable
{ 2. A szakdolgozat nyelvezete, stílusa, nyelvi helyessége. }
{ Nyelvezete mondatszerkesztése erősen kifogásolható, durva helyesírási hibákat tartalmaz. }
{ Nyelvezet, stílusa sok hiányossággal, kisebb helyesírási hibákat tartalmaz. }
{ Nyelvezete megfelelő, kevés stílushibával. }
{ Kifogástalan nyelvezet és stílus. }
{ 0 }
{  }

\EvaluationTable
{ 3. A dolgozat szakirodalmi hátterének feltárása, szakszerűsége, a hivatkozások helyessége. }
{ A szakirodalom nem releváns. A hivatkozások hiányoznak, vagy nem szakszerűek.  }
{ A szakirodalom feltárása szűk körű, csak tankönyvek, vagy jegyzetek anyagát tartalmazza. Hivatkozásai több helyen pontatlanok. }
{ Legfontosabb szakirodalmakat, korrekt hivatkozásokat tartalmaz a dolgozat. Az elmélet koherens, egységet alkot.  }
{ Legújabb széles körű szakirodalom alapján íródott a dolgozat. Hivatkozások pontosak. }
{ 0 }
{  }

\EvaluationTable
{ 4. A megvalósított feladat (program) nehézsége, összetettsége. }
{ A megvalósított feladat nem teljesíti a szakdolgozattól elvárt nehézségi szintet. }
{ A megvalósított feladat éppen teljesíti a szakdolgozattól elvárt nehézségi szintet.  }
{ A megvalósított feladat teljesíti a szakdolgozattól elvárt nehézségi szintet. }
{ A megvalósított feladat túlmutat a szakdolgozattól elvárt nehézségi szinten. }
{ 0 }
{  }

\EvaluationTable
{ 5. Megvalósítás minősége I.A feladat ismertetése, a célkitűzés definiálása (specifikáció, megvalósítási terv, összefüggések elemzése). }
{ A dolgozatból a feladat ismertetése, illetve a célkitűzés definiálása teljes mértékben hiányzik. }
{ A dolgozatban megjelenik a feladat ismertetése, továbbá a célkitűzés definiálása, de a leírás alapszintű, elnagyolt, hiányos. }
{ A dolgozatban a feladat ismertetése, illetve a célkitűzés definiálása formailag megfelelő, egy alapszintű leíráson túlmutat, szakmailag megfelelő, de nem teljes és ellentmondásmentes. }
{ A dolgozatban  a feladat ismertetése, illetve a célkitűzések definiálása szakmailag korrekt, formailag megfelelő, teljes és ellentmondásmentes. }
{ 0 }
{  }

\EvaluationTable
{ 6. Megvalósítás minősége II.Kódminőség. }
{ A forráskód tartalma, mennyisége, szerkezete nem megfelelő. }
{ A forráskód tartalma, mennyisége, szerkezete megfelelő, de nehézen áttekinthető, terjengős, strukturálatlan, és helyenként indokolatlan nyelvi elemeket tartalmaz.Nem tartalmazza a szükséges ellenőrzési, hibakezelési funkciókat. }
{ A forráskód tartalma, mennyisége, összetettsége jó. A forráskódban megjelennek a korszerű, hatékony nyelvi elemek. Jól strukturált, áttekinthető, tartalmazza a szükséges ellenőrzési, hibakezelési funkciókat. }
{ A forráskód tartalma, mennyisége, összetettsége kiváló. Korszerű, hatékony nyelvi elemeket tartalmaz. Jól strukturált, áttekinthető, öndokumentáló, újrafelhasználható. Megvalósított a teljeskörű hiba- és kivételkezelés. Erőforrás-ütemezés szempontjából optimalizált. }
{ 0 }
{  }

\EvaluationTable
{ 7. Megvalósítás minősége III.Fejlesztési- és felhasználói dokumentáció. }
{ A Fejlesztési- és felhasználói dokumentáció nem felel meg a követelményeknek.  }
{ A Fejlesztési- és felhasználói dokumentáció elnagyolt, hiányos, szerkezete logikátlan. }
{ A Fejlesztési- és felhasználói dokumentáció tartalmazza a fontosabb fejezeteket, szerkezete logikus, a fejezetek közötti kohézió jó. }
{ A Fejlesztési- és felhasználói dokumentáció kiváló színvonalú, a dokumentálási szabályoknak teljes mértékben megfelel. }
{ 0 }
{  }

\EvaluationTable
{ 8. Megvalósítás minősége IV.Tesztelés, futtatás. }
{ A dolgozat nem tartalmaz teszteléshez köthető részeket. A szoftver nem futtatható és/vagy nem a terveknek megfelelően működik. }
{ A dolgozatban megjelennek a teszteléssel kapcsolatos törekvések, de elnagyoltak, hiányosak. }
{ Jó színvonalú, de nem teljes körű tesztelés. }
{ Teljes körű és jól dokumentált validáció és verifikáció (formális módszerek, tesztelési terv, tesztelési jegyzőkönyv és jelentés). }
{ 0 }
{  }

\EvaluationTable
{ 9. A szakdolgozat összefoglalása. }
{ Zavaros, leíró összefoglalás, tézisek nélkül. }
{ Leíró jellegű összefoglalás, elnagyolt. }
{ Világos tagolt összefoglalás, korrekt. }
{ Kifogástalan, lényegre törő, továbbtervező. }
{ 0 }
{  }

\EvaluationTable
{ 10. Összbenyomás, konzulens/opponens véleménye }
{ A dolgozat szakmailag nem releváns, tartalmi formai követelményeknek nem felel meg. }
{ A dolgozat szakmailag kevésbé releváns, tartalmi formai követelményeknek megfelel. }
{ A dolgozat korrekt szakmai és módszertani felkészültséget tükröz. }
{ A dolgozat kiváló szakmai felkészültséget tükröz. }
{ 0 }
{  }

\bigskip

Amennyiben az értékelési szempontok közül valamelyik nulla pont, akkor az összes pontszám is nulla. Értékelésnél fél pontok is adhatók.

\bigskip

\begin{center}

\begin{tabular}{|l|c|}
\hline

\textbf{Összes pontszám} &
\textbf{0}

\\
\hline

\end{tabular}

\end{center}

\bigskip

\textbf{A dolgozat rövid szöveges értékelése (opcionális):}

\vspace{0.3cm}



\vspace{1cm}

\textbf{Javasolt érdemjegy megállapítása (kötelező)} Elégtelen (1)

\vspace{0.3cm}

A szakdolgozat javasolt érdemjegye az összesített pontszám alapján.

\bigskip

\begin{tabular}{ll}

\textbf{Elégtelen:} & 8--25 pont \\

\textbf{Elégséges:} & 26--32 pont \\

\textbf{Közepes:} & 33--38 pont \\

\textbf{Jó:} & 39--44 pont \\

\textbf{Jeles:} & 45--50 pont \\

\end{tabular}

\bigskip


\textbf{A bíráló javaslatai a védéshez (opcionális):}

\vspace{0.3cm}



\vspace{1cm}

\textbf{A hallgató által megválaszolandó kérdések (kötelező):}

\vspace{0.3cm}

1. sadas

2. sadasd

\vspace{1cm}

\begin{tabularx}{\textwidth}{lX}

\textbf{A szakdolgozat értékelése:} &
Elégtelen (1)

\end{tabularx}

\vspace{1.5cm}

\begin{tabularx}{\textwidth}{X X}

Eger, \today &
\hfill Konzulens aláírása

\end{tabularx}


\end{document}
